% Source-aligned English for printed p.123.
\SourcePage{123}

\LayerTarget{tohoku:section:1.1:p4}
\noindent bijective (the converse being false in general).

\LayerTarget{tohoku:section:1.1:p5}
The composite of two monomorphisms (epimorphisms) is a monomorphism
(epimorphism), hence the composite of two bijections is a bijection; likewise,
the composite of two isomorphisms is an isomorphism. If the composite $vu$ of
two morphisms $u,v$ is a monomorphism (respectively, an epimorphism), then $u$
(respectively, $v$) is also one. Although the development of such sorites is
obviously necessary, we shall most often dispense with spelling them out
explicitly in what follows, and shall content ourselves with indicating the
definitions carefully.

\LayerTarget{tohoku:section:1.1:p6}
Consider two monomorphisms $u:B\longrightarrow A$ and
$u':B'\longrightarrow A$. We say that \textit{$u'$ majorizes or contains $u$}
and write $u\leq u'$ if $u$ can be factored as $u'v$, where $v$ is a morphism
from $B$ to $B'$ (which is then uniquely determined). This is a
\textit{preorder} on the class of monomorphisms with codomain $A$. We say that
two such monomorphisms $u,u'$ are \textit{equivalent} if each majorizes the
other; then the corresponding morphisms $B\longrightarrow B'$ and
$B'\longrightarrow B$ are inverse to one another. Choose (for example, by
means of Hilbert's all-purpose symbol $\tau$) one monomorphism in each class of
equivalent monomorphisms: the chosen monomorphisms will be called the
\textit{subobjects of $A$}. Thus a subobject of $A$ is not a mere object of
$C$, but an object $B$ equipped with a monomorphism $u:B\longrightarrow A$,
called the \textit{canonical injection} of $B$ into $A$. (Nevertheless, by an
abuse of language, a subobject of $A$ will often be denoted by the symbol $B$
for the corresponding object $B\in C$.) The majorization relation defines an
\textit{order relation} (not merely a preorder) on the class of subobjects of
$A$. It follows from what was seen above that the subobjects of $A$ contained
in a subobject $B$ identify with the subobjects of $B$, this correspondence
respecting the natural order relations. (This does not, however, mean that a
subobject of $B$ is \textit{equal} to a subobject of $A$, which would indeed
require $A=B$.)

\LayerTarget{tohoku:section:1.1:p7}
Dually, consideration of a preorder on the class of epimorphisms from $A$
makes it possible to define the ordered class of \textit{quotient objects of
$A$}.

\LayerTarget{tohoku:section:1.1:p8}
Let $A\in C$, and let $(u_i)_{i\in I}$ be a nonempty family of morphisms
$u_i:A\longrightarrow A_i$. Then for every $B\in C$, the maps
$v\longrightarrow u_iv$ from $\mathrm{Hom}(B,A)$ to $\mathrm{Hom}(B,A_i)$
define a natural map
\[
  \mathrm{Hom}(B,A)\longrightarrow
  \prod_{i\in I}\mathrm{Hom}(B,A_i).
\]
We say that the $u_i$ define a \textit{representation of $A$ as a direct
product of the $A_i$} if, for every $B$, the preceding map is bijective. If
this is so, and if $A'$ is another object of $C$ represented as a product of
the $A'_i$ by morphisms $u'_i:A'\longrightarrow A'_i$ (the index set being the
same), then for every family $(v_i)$ of morphisms
$v_i:A_i\longrightarrow A'_i$ there exists a unique morphism $v$ from $A$ to
$A'$ such that $u'_iv=v_iu_i$ for every $i$. It follows that if the $v_i$ are
equivalences, then so is $v$; in particular, if the $v_i$ are the identity
maps $I_{A_i}$, we see that two objects $A,A'$ represented as products of the
family of $A_i$ are canonically isomorphic. It is therefore

\TohokuContinuation{124}
